\section{Data: incident streams, external events and online signals}\label{sec:data}

The framework of \cref{sec:framework} needs three kinds of input:
timestamped streams of recorded incidents, a calendar of salient external
events (the shock times $\tau_e$ and sizes $S_e$), and, where available,
binned series of online activity. This section reviews candidate sources
for each. For every source we record what can be checked in the provider's
own documentation: the timestamp the record carries, its temporal and
geographic resolution, the anti-Jewish bias category, and how the data can
be obtained. Properties that we could not confirm from primary documentation
at the time of writing (November 2026) are marked as such and are not used
in any argument. No real data have been analysed for this paper, and nothing
in this section is an empirical finding about antisemitism. Coverage,
definitional harmonisation and under-recording are treated in the companion
paper \citep{kornsteinlab2026darkfigure}; here the focus is on \emph{time}.

\subsection{What a point-process analysis needs from a record}\label{sec:data-needs}

Three clocks are attached to each incident~$i$ in stream~$k$: the
\emph{occurrence} time $t_i$ at which the incident happened; the
\emph{report} time $r_i = t_i + d_i$ at which it reached the recording body;
and the \emph{availability} time at which the record enters a file that an
analyst can obtain. The model of \cref{sec:framework} is defined on
occurrence times. Two departures from that ideal are common, and they bias
contagion estimates in different ways.

\paragraph{Report dates in place of occurrence dates.}
If the analysis uses $r_i$ instead of $t_i$, the observed process is the
true process with every point displaced by an independent delay $d_i \sim
F_k$. The number of points is preserved, but for distinct points the second-order
structure is convolved with the law of $d_i - d_j$: a parent--offspring pair
at lag $u$ appears at lag $u + d_j - d_i$, which is negative with positive
probability whenever delays are dispersed. The displaced process need not be
a Hawkes process. Contagion estimated from report dates therefore mixes the
excitation kernel with the delay distribution, and this matters most when
typical delays are comparable with the kernel time scale $1/\beta$.

\paragraph{Coarse time stamps.}
Many sources record dates without times, and some record only the month of
occurrence. Tied time stamps have probability zero under a
continuous-time model. They can be handled by randomising times within the recorded
interval and averaging over randomisations, or by working with the
discrete-time INAR representation of \cref{eq:inar} \citep{kirchner2016},
which is the natural choice for daily counts. A stream whose occurrence time
is known only to the month cannot identify kernels on the scale of days
from occurrence times alone.

\subsection{Incident streams}\label{sec:data-incidents}

\paragraph{City-level open data: NYPD hate crimes.}
The New York City Police Department publishes a file of confirmed hate-crime
incidents on the NYC Open Data portal (dataset \texttt{bqiq-cu78})
\citep{nypdhate}, beginning in 2019. According to the portal's data
dictionary, \emph{Complaint Year Number} and \emph{Month Number} give the
year and month in which the incident occurred, and \emph{Record Create Date}
gives the date on which the report was filed. A \emph{Bias Motive
Description} field holds the NYPD's bias category, whose values include
\texttt{ANTI-JEWISH}. Geography is given by the complaint precinct, the
patrol borough and the county; each of the city's five counties is
coterminous with a borough. For point-process work the file thus has a
two-level time structure: the occurrence time is interval-censored to a
calendar month, while the report date is known to the day. We treat the
record-creation date as a report date and model the delay between occurrence
and record creation explicitly (\cref{sec:data-delays}), with the
occurrence month as a constraint. Whether day-level occurrence dates can be
obtained from the NYPD, for example under a data-use agreement, is an open
feasibility question for the pilot (\cref{sec:plan}); we do not assume it.
Changes to reporting systems are structural breaks for time-series use. The
New York State Comptroller reported that the state's Division of Criminal
Justice Services did not publish 2023 hate-crime data for New York City
because of the NYPD's transition in data-reporting systems as part of its
move to NIBRS \citep{osc2024hate}. \Citet{rosenbaum2025} analyses NYPD
administrative hate-crime records for 2019--2024 at monthly resolution.
Adjusting for the Jewish population of each borough, the study reports
that monthly anti-Jewish hate crimes were on average about twice as common
in the first year of the Israel--Hamas war as in the preceding years of the
study period (prevalence ratio 1.97, 95\% CI 1.64--2.35). Monthly
resolution suffices for such before--after contrasts but not for kernels
on the scale of days.

\paragraph{FBI NIBRS.}
The FBI's Uniform Crime Reporting Program retired its Summary Reporting
System on 1 January 2021 and moved to collection through the National
Incident-Based Reporting System (NIBRS) \citep{fbi2026nibrs}. For hate
crimes, however, the FBI reopened the 2021 collection to summary-format
submissions \citep{fbi2023supp}. According to the Congressional Research
Service, it has accepted both formats since 2022 \citep{crs2025hate}. Each
NIBRS incident carries the mandatory Data Element~3, \emph{Incident
Date/Hour}. The date is recorded as \texttt{YYYYMMDD} and the hour on a
24-hour clock. When the incident date is unknown, the report date is entered
instead and flagged with a report-date indicator
\citep{fbi2025nibrs,fbi2025nibrstech}. Anti-Jewish bias is code~21 of Data
Element~8A, \emph{Bias Motivation} \citep{fbi2025nibrs}. The geographic unit
is the reporting agency, identified by its ORI number. NIBRS submissions
thus supply occurrence dates, and hours where recorded, for individual
anti-Jewish incidents. For summary-format submissions we have not confirmed
comparable date fields, so the stream is restricted to NIBRS records. The
data are also slow. National statistics for data year 2024 were released on
5 August 2025 \citep{fbi2025release}, so for that year the lag from
occurrence to national availability was roughly 7 to 19 months, depending
on when in the year an incident occurred. Agency participation varies by
agency and year; the companion paper models it.

\paragraph{ADL Audit and H.E.A.T.\ Map.}
ADL describes its annual Audit as covering criminal and non-criminal
incidents of harassment, vandalism and assault, as reported to ADL by
victims, law enforcement, the media and partner organisations
\citep{adl2025audit}. An ADL stream therefore measures a different quantity
from a police stream, and the two reach the record through different funnels
(\cref{sec:data-thinning}). ADL states that since 7 October 2023 the Audit
has included certain expressions of opposition to Zionism and of support for
violent resistance against Israel or Zionists, while continuing to exclude
legitimate political protest \citep{adl2025audit}. For a point-process model
this matters for a structural reason, separate from any substantive
judgement. A change in recording scope at a date that is also a candidate
shock time $\tau_e$ is not identified separately from the shock term unless
the change is modelled explicitly. An ADL stream therefore needs a scope-break
indicator in $\mu_k(t)$, or records coded under a constant scope rule. ADL's
H.E.A.T.\ Map (Hate, Extremism, Antisemitism, Terrorism) publishes
incident-level listings that can be downloaded as a CSV file, with some
details removed to protect victims' privacy \citep{adlheatmap}. We have not
confirmed whether the listed date is the date of occurrence or of report to
ADL, nor how redaction affects dates. These checks are part of the pilot's
feasibility stage.

\paragraph{CST (United Kingdom).}
The Community Security Trust records antisemitic incidents in the United
Kingdom and publishes annual reports, most recently for 2025
\citep{cst2026incidents}. CST data would serve as an international comparison
stream. That report illustrates a further threat to time-series analysis.
CST states that since March 2024 a technical problem has disrupted the flow
of reports to it from the Metropolitan Police \citep{cst2026incidents}. In
the notation of \cref{sec:data-thinning}, an interruption in a referral feed
is a step change in the recording probability $p_k(t)$, which a model with
constant $p_k$ would read as a change in the incident process itself.
Incident-level records with dates would require a data-sharing agreement,
and a CST stream is therefore conditional.

\subsection{Reporting delays and nowcasting}\label{sec:data-delays}

Three lags separate an incident from an analyst's file: occurrence to report
($d_i$, driven by victims' and witnesses' decisions); report to
classification or confirmation as a bias incident; and confirmation to
publication, which is set by each provider's release schedule. Suppose that
recording and delays are independent of the history and across incidents,
with recording probability $p_k$ (\cref{sec:data-thinning}) and total-delay distribution
$F_k$. Then the occurrence times of the stream-$k$ incidents that are
available at analysis time $T_a$ form an independent thinning of the true
process, with time-varying retention probability $p_k F_k(T_a - t)$. Hence
\begin{equation}\label{eq:data-righttrunc}
  \E\bigl[N_k^{\mathrm{avail}}(\dd t)\bigr]
  = p_k\,F_k(T_a - t)\,\E\bigl[\lambda_k(t)\bigr]\,\dd t .
\end{equation}
Near $T_a$ the factor $F_k(T_a - t)$ is small. Without correction, a surge
that is still under way looks as if it is already subsiding, which is
exactly the failure an alert system must avoid. The correction is a
nowcast, as developed for infectious-disease surveillance
\citep{hohle2014,mcgough2020} and adapted to the Hawkes setting in
\cref{prop:nowcast}. Estimating $F_k$ needs records that carry both an
occurrence and a report time. In the NYPD file the month of occurrence and
the record-creation date bound each delay. NIBRS supplies the occurrence date
directly, but its release lag makes it a source for retrospective estimation
of $\Amat$ rather than for alerting. Open-data files may also be revised after
first release, so every analysis will be tied to an archived, checksummed
vintage of each file.

\subsection{External-event calendars}\label{sec:data-events}

The shock term $\sum_e \beta_k S_e \psi_k(t - \tau_e)$ requires event times
and sizes that are fixed independently of the incident data. Two established
event datasets can supply candidate events: ACLED \citep{raleigh2010} and
GDELT \citep{leetaru2013}. We have not re-verified their current
geographic and temporal coverage or their licence terms for this draft.
Before registration we will record the terms in force on the download date
and comply with them, including any limits on redistributing raw records.
Where redistribution is not permitted, the replication package will contain
event identifiers, dates and derived salience scores together with
re-download scripts, not the raw records.

Neither dataset settles which events are \emph{salient} for the outcome. The
risk is that the analyst picks events after seeing surges, which would bias
the shock loadings $\beta_k$ upwards and the contagion terms downwards. We
therefore fix the event list before outcome data are examined, using a
content-neutral procedure set out in full in \cref{sec:plan-events}:
\begin{enumerate}[(i)]
  \item event classes are defined in advance by type (for example, armed
  escalations, mass-casualty attacks, high-profile public statements),
  symmetrically and without reference to the parties involved;
  \item salience is measured by a pre-specified, outcome-independent index,
  such as the volume of news coverage relative to its own trailing baseline,
  and an event enters the list when the index crosses a fixed threshold;
  \item coding is done by coders who have no access to incident data; and
  \item the list of times $\tau_e$, classes and sizes $S_e$ is frozen and
  time-stamped before incident data for the study period are loaded.
\end{enumerate}
Calendar effects such as religious holidays and school terms are recurring
and predictable, so they enter the background $\mu_k(t)$, not the shock
list.

\subsection{Online signals}\label{sec:data-online}

\paragraph{Platform data access.}
Article~40 of the EU Digital Services Act provides a procedure by which
vetted researchers may obtain data from very large online platforms and
very large online search engines for research that contributes to the detection,
identification and understanding of systemic risks in the Union. It also
contains a separate provision on access to publicly accessible data, and it
leaves technical conditions to delegated acts \citep{eu2022dsa}. Whether a
given project qualifies, and on what terms, is decided through that
procedure. We have not confirmed its operational status for the platforms
and periods of interest, and the plan in \cref{sec:plan} treats online
streams as optional, so that the offline analysis does not depend on them.
The terms of any research interface used will be recorded at registration.

\paragraph{Classifier-based volume series and their measurement error.}
An online stream enters the model as a binned count series, which is handled
by the INAR approximation of \cref{eq:inar} or a binned Hawkes likelihood.
It is constructed by applying a classifier to a defined corpus, following
the lab's proposals on detection and measurement
\citep{kornsteinlab2026detection,kornsteinlab2026measurement}. Two points
from those proposals carry over. First, the estimand concerns the sampled
corpus, not platform-wide prevalence. Second, classifier error must be
reported through both precision and recall, not a single score. For the
time-series use, let $U_t$ be the number of posts in bin $t$, $V_t$ the
number that are antisemitic under the reference annotation, and
$\widehat V_t$ the number flagged. If each post is flagged independently
with probability $\eta_1$ when antisemitic and $\eta_0$ otherwise, then
\begin{equation}\label{eq:data-classerr}
  \E\bigl[\widehat V_t \mid U_t, V_t\bigr]
  = \eta_1 V_t + \eta_0\,(U_t - V_t),
  \qquad
  \widetilde V_t = \frac{\widehat V_t - \eta_0 U_t}{\eta_1 - \eta_0}
  \quad (\eta_1 > \eta_0),
\end{equation}
where $\widetilde V_t$ is conditionally unbiased for $V_t$ when $\eta_0$ and
$\eta_1$ are known. The uncorrected series loads on total volume $U_t$ with
weight $\eta_0$. If total posting volume rises after a salient external
event, the uncorrected series rises even when $V_t$ does not, and a common
shock can then masquerade as online-to-offline excitation. The error rates
are not fixed constants. The detection proposal singles out quotation and
criticism of antisemitic expression as hard cases, and if the share of such
posts changes after events, $\eta_0$ moves with the shocks. The plan
therefore estimates $(\eta_0, \eta_1)$ on validation samples stratified by
period, propagates their uncertainty, and reports corrected and uncorrected
analyses (\cref{sec:plan-sens}). Any estimated online-to-offline coefficient
$\alpha_{\mathrm{off},\mathrm{on}}$ is interpreted as Granger-type
predictive dependence in the sense of \citet{eichler2017}
(\cref{prop:granger}), not as a structural causal effect.

\subsection{Prior empirical literature}\label{sec:data-lit}

\Citet{mohler2011} developed self-exciting point-process models of
crime; \citet{reinhart2018} reviews spatio-temporal versions and their
applications. Related applications include a spatiotemporal test of whether
gun violence is contagious \citep{loeffler2018} and self-exciting models of
civilian deaths in Iraq \citep{lewis2012}. For hate crime specifically,
\citet{king2013} study the temporal clustering of hate-motivated offending.
Two studies examine links between online content and offline hate crime:
\citet{muller2021} study social media and hate crime, and
\citet{williams2020} examine anti-Black and anti-Muslim social-media posts
as predictors of offline racially and religiously aggravated crime. For
anti-Jewish incidents, \citet{rosenbaum2025} compares monthly NYPD counts
before and after a salient external event (see above). We use this
literature to motivate the questions, not as prior estimates. None of these
studies estimates a multivariate branching matrix for antisemitic incidents,
and our framework does not take its parameter values from them. Separating
exogenous shocks from contagion is the identification problem of
\cref{sec:identification}.

\subsection{How the dark-figure estimates enter: recorded streams are thinned}\label{sec:data-thinning}

Recorded incidents are a subset of true incidents. Under the working
assumption that each true incident in stream $k$ is recorded independently
with probability $p_k$, the recorded stream is a $p_k$-thinning of the true
one. By \cref{thm:thinning}, contagion estimated from recorded data then
targets an attenuated apparent branching ratio, not the true one. The
companion paper supplies $p_k$. In its notation the dark-figure ratio
relative to list $k$ is $D_k = N/n_{k+}$, so $p_k$ is estimated by
$1/\widehat D_k$. Its funnel
decomposition shows that $p_k$ depends on the source: a city police stream
such as the NYPD file is reached through disclosure, reporting to police and
bias classification ($p \approx \rvic\rpol\rcls$); the FBI stream adds
agency submission ($p \approx \rvic\rpol\rcls\rsub$); and an NGO audit has a
funnel of its own \citep{kornsteinlab2026darkfigure}. Two cautions follow.
First, the companion paper's estimates refer to coarse domains, such as
state-by-year cells, while
contagion operates on the scale of days within a city. Using them assumes
that $p_k$ is constant within a domain, and this assumption is varied in the
sensitivity analysis. Second, recording may itself respond to salient
events, for instance if heightened attention raises the propensity to report.
Recording would then no longer be independent of the history, the recorded
surge would overstate the true surge, and \cref{thm:thinning} would not
apply as stated. The plan treats time-varying $p_k$ as a pre-specified
sensitivity analysis (\cref{sec:plan-sens}).

\Cref{tab:data-streams} summarises the candidate streams.

\begin{table}[tbp]
\centering
\caption{Candidate data streams for the pre-registered analysis. ``TBC'':
not confirmed from the provider's documentation at the time of writing; to
be confirmed before registration. Delays are from occurrence to
availability to analysts.}
\label{tab:data-streams}
\footnotesize
\setlength{\tabcolsep}{3pt}
\renewcommand{\arraystretch}{1.15}
\begin{tabularx}{\textwidth}{@{}>{\raggedright\arraybackslash}p{2.35cm}
  >{\raggedright\arraybackslash}X
  >{\raggedright\arraybackslash}p{2.2cm}
  >{\raggedright\arraybackslash}X
  >{\raggedright\arraybackslash}p{2.6cm}@{}}
\toprule
Stream & Time stamp (resolution) & Geography & Delay to availability & Access \\
\midrule
NYPD hate crimes \citep{nypdhate}
  & Occurrence: month; record creation: day. From 2019.
  & Precinct, patrol borough, county
  & Report delay bounded by occurrence month and record date; refresh schedule TBC
  & Public download (NYC Open Data) \\
FBI NIBRS \citep{fbi2025nibrs,fbi2026nibrs,crs2025hate}
  & Incident date and hour; report date substituted (flagged) if incident date unknown. NIBRS from 2021; hate-crime collection also accepts summary format since 2022.
  & Reporting agency (ORI)
  & Annual national release; data year 2024 released 5 Aug.\ 2025
  & Public (FBI Crime Data Explorer) \\
ADL Audit / H.E.A.T.\ Map \citep{adl2025audit,adlheatmap}
  & Incident date listed; occurrence vs.\ report date TBC. Scope change from 7 Oct.\ 2023.
  & Place listed; resolution TBC
  & TBC
  & Public map; CSV download with redactions \\
CST (UK) \citep{cst2026incidents}
  & TBC
  & TBC
  & Annual report; 2025 report published Feb.\ 2026
  & Incident-level records by agreement (TBC) \\
ACLED \citep{raleigh2010}
  & TBC
  & Event location
  & TBC
  & Provider terms (TBC) \\
GDELT \citep{leetaru2013}
  & TBC
  & Event location
  & TBC
  & Provider terms (TBC) \\
Platform data \citep{eu2022dsa}
  & Post time stamps (binned for analysis)
  & Varies
  & Depends on access route
  & DSA Art.~40 procedure; status TBC \\
Classifier volume \citep{kornsteinlab2026detection,kornsteinlab2026measurement}
  & Bin chosen by design (e.g.\ day)
  & None or coarse
  & Construction lag only
  & Lab-constructed from a defined corpus \\
\bottomrule
\end{tabularx}
\end{table}
